% ============================================================================
% GÉNÉRÉ par scripts/exemples-guide.ts — NE PAS ÉDITER À LA MAIN.
% Régénérer : npm run exemples (chaque chiffre est vérifié contre le moteur).
% ============================================================================
\pexemples Le montant maximal (par enfant, plus le supplément monoparental le cas
échéant) est réduit de \pc{4} du $\RFN$ excédant le seuil,
sans jamais descendre sous le montant minimal. Le supplément pour fournitures
scolaires (\D{124} par enfant de 4 à 16~ans) s'ajoute
sans égard au revenu.

\medskip\noindent\textbf{M6 --- famille monoparentale, 35~ans, revenu de travail de \D{35000}, un enfant : 3~ans (garde subventionnée, \D{2000}).}
$\RFN$ de \num{33265}, sous le seuil monoparental de \num{43280} :
aucune réduction, l'allocation est maximale.
\begin{align*}
\text{maximum} &= \num{3006} + \num{1055} = \num{4061} \\
\text{allocation} &= \num{4061} - 0 = \num{4061}
\end{align*}
L'enfant a 3~ans : pas de supplément pour fournitures scolaires
(\D{0} — il vise les 4 à 16~ans).

\medskip\noindent\textbf{M8 --- couple, 38 et 36~ans, revenus de travail de \D{60000} et \D{40000}, deux enfants : 4~ans (garde non subventionnée, \D{13000}) et 8~ans (garde non subventionnée, \D{3000}).}
\begin{align*}
\text{maximum} &= 2 \times \num{3006} = \num{6012} \\
\text{réduction} &= \pos{\num{96230} - \num{59369}} \times \pc{4} = \num{1474.44} \\
\text{allocation} &= \max\bigl(\num{2392},\ \num{6012} - \num{1474.44}\bigr) = \num{4537.56}
\end{align*}
Le plancher (montant minimal, \num{2392}) ne mord pas encore :
allocation de \D{4537.56}, plus \D{248} de fournitures scolaires
($2 \times \num{124}$, enfants de 4 et 8~ans).

\medskip\noindent\textbf{M9 --- couple, 45 et 45~ans, revenus de travail de \D{120000} et \D{60000}, un enfant : 5~ans (garde non subventionnée, \D{15000}).}
$\RFN$ de \num{175521} : la réduction (\num{4646.08}) dépasse
l'écart entre maximum (\num{3006}) et minimum (\num{1196}) —
l'allocation est au \emph{plancher}.
\begin{align*}
\text{allocation} &= \max\bigl(\num{1196},\ \num{3006} - \num{4646.08}\bigr) = \num{1196}
\end{align*}
Toute famille admissible reçoit au moins le minimum, ici \D{1196},
plus \D{124} de fournitures scolaires.
